% Part: sets-functions-relations
% Chapter: arithmetization
% Section: cuts
%
\documentclass[../../../include/open-logic-section]{subfiles}

\begin{document}

\olfileid[tr]{sfr}{arith}{cuts}
\olsection{$\Rat$'den $\Real$'e}

Tamlık Özelliği özünde, gerçel doğru üzerindeki her $\alpha$ noktasının bu
doğruyu kusursuz biçimde iki yarıya ayırdığını gösterir: $\alpha$'nın en küçük
üst sınır olduğu yarı ve $\alpha$'nın en büyük alt sınır olduğu yarı. Gerçel
sayıları rasyonel sayılardan \emph{kurmak} için Dedekind, gerçel sayıları
rasyonel sayıları bölen \emph{kesimler} olarak düşünmemizi önerdi. Yani
$\sqrt{2}$'yi, $< \sqrt{2}$ olan rasyonel sayıları $> \sqrt{2}$ olan
rasyonel sayılardan ayıran \emph{kesim} ile özdeşleştiririz.

Bunu biraz düzenleyelim. Rasyonel sayıları iki yarıya bölersek, yaptığımız
bölmeyi yalnızca \emph{alt} yarısını ele alarak tekil biçimde belirleyebiliriz.
Dolayısıyla kesin bir tanım verelim:

\begin{defn}[Kesim]
Bir \emph{kesim} $\alpha$, rasyonel sayıların en büyük elemanı olmayan, boş
olmayan herhangi bir öz başlangıç kesitidir. Yani $\alpha$ ancak ve ancak şu
koşullarda bir kesimdir:
\begin{enumerate}
	\item \emph{boş olmayan ve öz}: $\emptyset \neq \alpha \subsetneq \Rat$
	\item \emph{başlangıç kesiti}: bütün $p,q \in \Rat$ için,
	$p < q \land q \in \alpha$ ise $p \in \alpha$
	\item \emph{en büyük elemanı yok}: her $p \in \alpha$ için,
	$p < q$ olacak biçimde bir $q \in \alpha$ vardır
\end{enumerate}
O hâlde $\Real$, kesimler kümesidir.
\end{defn}

Artık $\sqrt{2} = \Setabs{p \in \Rat}{p^2 < 2\text{ veya }p < 0}$
diyebiliriz. Elbette bunun gerçekten bir kesim olduğunu denetlemeliyiz; fakat
bu denetimi \olref[check]{sec}'e bırakıyoruz.

Daha önce olduğu gibi, bazı varlıkları tanımladıktan sonra onlar üzerindeki
temel fonksiyon ve bağıntıları da tanımlamamız gerekir. Kolay bir tanımla
başlayalım:
\begin{align*}
	\alpha \leq \beta \text{ ancak ve ancak }\alpha \subseteq \beta
\end{align*}
Bu sıralama tanımı, kesimler kümesinin Tamlık Özelliği'ne sahip olduğu
yönündeki temel sonucu \emph{ifade etmemizi} sağlar. Tam olarak açarsak önerme
şu biçimdedir: $S$, üstten sınırlı ve boş olmayan bir kesimler kümesiyse
$S$'nin en küçük bir üst sınırı vardır. Daha ayrıntılı olarak, $S$ için üst
sınır olan, yani $(\forall \alpha \in S)\alpha \subseteq \lambda$ koşulunu
sağlayan bir $\lambda$ kesimi vardır ve $\lambda$ böyle kesimlerin en
küçüğüdür; yani
$(\forall \beta \in \Real)((\forall \alpha \in S)\alpha \subseteq \beta
\lif \lambda \subseteq \beta)$. Şimdi sonucun ispatını verelim:

\begin{thm}\ollabel{realcompleteness}
Kesimler kümesi Tamlık Özelliği'ne sahiptir.
\end{thm}

\begin{proof}
$S$, üstten sınırlı ve boş olmayan herhangi bir kesimler kümesi olsun.
$\lambda = \bigcup S$ olsun.
%\Setabs{p \in \Rat}{(\exists \alpha \in S)p \in \alpha}$$
Önce $\lambda$'nın bir kesim olduğunu ileri sürüyoruz:
\begin{enumerate}
\item $S$ boş olmadığından en az bir kesim $S$'dedir; dolayısıyla
$\emptyset \neq \lambda$. $S$ bir kesimler kümesi olduğundan
$\lambda \subseteq \Rat$. $S$'nin bir üst sınırı bulunduğundan, bazı
$p \in \Rat$ her $\alpha \in S$ kesiminde bulunmaz. Öyleyse
$p\notin \lambda$ ve dolayısıyla $\lambda \subsetneq \Rat$.
\item $p < q \in \lambda$ olduğunu varsayın. O hâlde $q \in \alpha$ olacak
biçimde bir $\alpha \in S$ vardır. $\alpha$ bir kesim olduğundan
$p \in \alpha$; dolayısıyla $p \in \lambda$.
\item $p \in \lambda$ olduğunu varsayın. O hâlde $p \in \alpha$ olacak
biçimde bir $\alpha \in S$ vardır. $\alpha$ bir kesim olduğundan
$p < q$ olacak biçimde bir $q \in \alpha$ vardır. Dolayısıyla
$q \in \lambda$.
\end{enumerate}
Bu, ileri sürülen iddiayı ispatlar. Dahası,
$(\forall \alpha \in S)\alpha \subseteq \bigcup S = \lambda$ olduğu açıktır;
yani $\lambda$, $S$ için bir üst sınırdır. Şimdi $\beta \in \Real$'in
de bir üst sınır olduğunu, yani
$(\forall \alpha \in S)\alpha \subseteq \beta$ olduğunu varsayın. Her
$p \in \Rat$ için, $p \in \lambda$ ise $p \in \alpha$ olacak biçimde bir
$\alpha \in S$ vardır; böylece $p \in \beta$. Genelleştirirsek
$\lambda \subseteq \beta$ olur. Demek ki $\lambda$, $S$'nin \emph{en küçük}
üst sınırıdır.
\end{proof}

Böylece Tamlık Özelliği'ni sağlayan bir grup varlığımız oldu. Bunu söylemenin
bir yolu da şudur: kesimlerimizde “boşluk” yoktur. (Dolayısıyla rasyonel
sayılar yerine gerçel sayılardan yeni “kesimler” almak, ilginç yeni nesneler
vermez.)

Sonra gerçel sayılar üzerinde bazı işlemler tanımlamalıyız. Her $p \in \Rat$
için $p_\Real = \Setabs{q \in \Rat}{q < p}$ olduğunu belirleyerek rasyonel
sayıları gerçel sayılara gömmekle başlarız. Ardından şunları tanımlarız:
\begin{align*}
	\alpha + \beta &= \Setabs{p + q}{p \in \alpha \land q \in \beta}\\
%	\alpha - \beta &\defis \Setabs{p - q}{p \in \alpha \land q \in \Rat \setminus \beta}\\
	\alpha \times \beta &=
	\Setabs{p \times q}{0 \leq p \in \alpha \land 0 \leq q \in \beta} \cup 0_\Real & \text{eğer }\alpha, \beta \geq 0_\Real
%	\alpha \div \beta &\defis
%	\Setabs{\nicefrac{p}{q}}{p \in \alpha \land q \in \Rat \setminus \beta}  & \text{if }\alpha \geq 0_\Real, \beta > 0_\Real
\end{align*}
Diğer çarpma durumlarını ele almak için önce şunu tanımlayalım:
%that $0_\Real \times \alpha = 0_\Real = \alpha \times 0_\Real$, and add:
\begin{align*}
	-\alpha &= \Setabs{p - q}{p,q \in \Rat \land p < 0 \land q \notin \alpha}
\end{align*}
ve ardından şunları belirleyelim:
\begin{align*}
	\alpha \times \beta &\defis
	\begin{cases}
		0_\Real &\text{eğer }\alpha = 0_\Real\text{ veya }\beta = 0_\Real\\
		\mathord{-}\alpha \times \mathord{-}\beta &\text{eğer }\alpha < 0_\Real\text{ ve }\beta < 0_\Real\\
		\mathord{-}(\mathord{-}\alpha \times \beta) &\text{eğer }\alpha < 0_\Real \text{ ve }\beta > 0_\Real\\
		\mathord{-}(\alpha \times \mathord{-}\beta) &\text{eğer }\alpha > 0_\Real \text{ ve }\beta < 0_\Real
	\end{cases}
\end{align*}
Şimdi bu tanımların her birinin daima bir kesim verdiğini denetlemeliyiz. Son
olarak bu biçimde tanımlanan kesimlerin tam da gerektiği gibi davrandığını
kolay ama uzun bir gösterimle ortaya koymalıyız. Ancak bunların hepsini
\olref[check]{sec}'e bırakıyoruz.
\end{document}
